\section{Scope, source audit, and principal conclusions}
\label{sec:context}

\subsection{The question addressed}

The practical question is whether repeated observations of a partially computed
knot complex can become sufficiently cheap to justify earlier recognition.
The asymptotic question is whether such observations can help the program avoid
the exponentially large intermediate representations that obstruct a general
quasi-polynomial bound. These questions require different evidence. A fast
determinant subroutine answers the first only on workloads where it is used
often enough. It answers the second only when accompanied by control of the
number and size of the complexes that must be constructed before it succeeds.

The source examined here is the public ProveIt revision
\begin{center}\basecommit.\end{center}
All references to the ``maintained'' implementation mean that pinned revision,
before this package's additions. The repository already contains 28 numbered
reports, numerous integrated improvements, and a Python implementation far
beyond the earliest scanning baseline. Its relevant files are
\path{fast/fastunknot/shadow_scan.py},
\path{fast/fastunknot/integer_determinant.py},
\path{fast/fastunknot/tait_blocks.py}, and the research prototype
\path{fast/determinant_research/boundary_tait.py}.
The maintained explanations are
\path{synthesis/determinant_continuations.tex},
\path{synthesis/boundary_reuse.tex},
\path{synthesis/tait_blocks.tex}, and
\path{synthesis/shadow_fallback.tex} \cite{proveit}.

Report 26 already supplies the marked four-residue obstruction and a
singular-safe rational terminal kernel. Later maintained work supplies cut-face
fragments, a boundary-only geometric acceptance check, articulation factorization
of direct cofactors, and a budgeted Euler fallback. We use those results openly.
Reintroducing a Schur complement or an Euler norm without this attribution would
misstate the continuation's contribution.

\subsection{What this continuation adds}

The deliverable has four connected parts.

\begin{enumerate}
\item An exact terminal count, $t=b/2$, and the resulting $b-2$ residual
dimension bound for a nonempty frontier. This sharpens the earlier loose linear
bound for the actual cut-face construction.
\item A modular observer with a fully proved integer interpretation. It combines
congruences by the Chinese remainder theorem (CRT), retains the exact coordinate
sum, and solves a four-dimensional integer norm minimization explicitly.
\item A deterministic threshold completion theorem. The modulus required to
decide whether the existing bound exceeds one is smaller than the modulus
required to reconstruct all its integer coordinates.
\item An optional implementation in the maintained package, with independent replay of its
positive observations and paired comparisons against both relevant predecessors.
\end{enumerate}

The basic modular-determinant idea is established numerical algebra
\cite{dumasurbanska}. Boundary response matrices also have a substantial
literature \cite{kenyonwilson}. Our claims concern the explicit proofs and
implementation given here; a limited literature search is not a certification
of historical priority for every elementary lemma.

\begin{samepage}
\begin{theorem}[Combined observer guarantee]
\label{thm:main}
Suppose a validated classical knot diagram has a proper scan prefix whose
relative complex has been obtained by the algebraically valid operations
implemented by the maintained scanner. Let $\mathcal C_a$ be its whole differential components, with positive
multiplicities $m_a$, and let $x_a\in\Z^4$ be the marked completed Euler vectors.
For a product $M$ of distinct validated odd primes, the algorithm below computes
a number $L_M$ satisfying
\[
 0\le L_M\le B_4:=\sum_a m_a\norm{x_a}
 \le \dim_{\F_2}\Khred(K;\F_2).
\]
It may therefore return \code{KNOTTED} whenever $L_M>1$.
If $|x_{a,r}|\le C_a$ and $M>\max_a C_a$, then
\[
             L_M>1\quad\Longleftrightarrow\quad B_4>1.
\]
For a suffix with $b>0$ frontier darts, each modular determinant query after
preparation uses a matrix of dimension at most $b-2$, or returns the proved
zero residue without forming that matrix. These statements remain valid at
primes where the interior matrix becomes more singular.
\end{theorem}
\end{samepage}

The proof is distributed through \cref{sec:euler,sec:geometry,sec:modular}.
The hypothesis about the relative complex is essential: a randomly chosen
matrix satisfying $d^2=0$ is not enough. Likewise, $L_M\le1$ is not an unknot
certificate, even when the threshold comparison with $B_4$ is complete.

\subsection{Current literature and the complexity target}

Lackenby's July 2026 preprint develops an explicit hierarchy framework and
efficiently encoded certificates \cite{lackenby2026}. Proposition 9.1 bounds
the number of iterations by
\[
                         L(g+1)^L,
\]
where $L$ is a hierarchy-length bound and $g$ a pattern-complexity bound.
The refined algorithm controls hierarchy length using handle complexity;
Theorem 14.1 establishes, under its manifold and handle-structure
hypotheses, the existence of an essential hierarchy with a polynomial-size
encoding and polynomial-time verification. The paper leaves the proposed speed-up to further work. These are
substantial algorithmic foundations, but polynomial certificate verification
does not supply a polynomial or quasi-polynomial certificate-construction time.

The exact exponent also deserves care. The March 2021 Oxford technical
handout states $2^{O((\log N)^3)}=N^{O((\log N)^2)}$ \cite{lackenbytalk}.
The Oxford seminar abstract instead states $N^{O(\log N)}$
\cite{lackenbyabstract}. Both are
quasi-polynomial, but the latter is stronger. We keep the repository's stronger
target visible while using $2^{(\log N)^{O(1)}}$ when discussing the broad class.
No result here converts one announced exponent into another.

A separate September 2026 preprint by Musick explicitly claims polynomial
unknot recognition \cite{musick}. This changes what a current literature note
must acknowledge. The present work does not validate or integrate that claimed
algorithm. A limited review identifies the grammar-to-geometric-local-minimum
step and total compressed-description accounting as obligations for a separate
audit; it supplies no counterexample or disproof. Our definite complexity claim
is about the maintained ProveIt implementation: it still has no established
general quasi-polynomial guarantee.

The representation obstacle is concrete. In their January 2026 preprint,
Kelom\"aki and Sch\"utz prove
exponential behavior for basis-oriented scanning and divide-and-conquer on
particular positive or alternating three-braids, while also obtaining
polynomial-time structural formulas for three-braid Khovanov homology
\cite{kelomaki}. Thus a small frontier and a large explicit homology basis can
coexist. A succinct answer may nevertheless be tractable. This continuation
compresses an observation of the complex; it does not prove that the complex
itself remains succinct.

\subsection{Reading the empirical conclusion}

On the favorable repeated-query workload, 47 interior vertices are eliminated
once and 200 actual boundary matchings are observed. The modular path is
$26.67$ times faster than the maintained direct path and $3.97$ times faster
than the rational prototype. At 127 interior vertices and 50 queries, the
corresponding ratios are $10.74$ and $10.22$.

Conversely, the all-terminal workload slows down, the smallest nontrivial
interior favors the rational prototype, and the raw recognizer examples all
favor the existing direct observer. These measurements justify an optional
backend and further scheduling research. They do not justify changing the
default recognizer or claiming a general recognition speed-up.

\begin{center}
\begin{tabularx}{\textwidth}{@{}P{0.28\textwidth}X@{}}
\toprule
Claim category & Status in this package\\
\midrule
Modular arithmetic soundness & Proved, implemented, independently audited.\\
Exact terminal count & Proved for the checked cut-face representation.\\
Deterministic threshold completion & Proved; implementation reports whether its finite palette suffices.\\
Repeated-query acceleration & Measured with setup included and both predecessors compared.\\
Whole-recognizer acceleration & Not demonstrated on the tested corpus.\\
General quasi-polynomial recognition & Not established by this work.\\
\bottomrule
\end{tabularx}
\end{center}
